\section{问题重述}

\subsection{问题背景}

近年来，大语言模型（Large Language Models，LLMs）在知识问答、代码生成和复杂推理等任务中表现出显著的规模效应。经典标度律表明，验证损失通常随模型参数量和训练 Token 数增加而下降，这为估计不同训练规模下的模型性能提供了经验基础。然而，参数量和数据量只描述了训练资源的数量维度：相同规模的训练数据可能具有不同质量，不同领域的组合方式也会改变模型在各类验证任务上的表现。此外，预训练 Loss 与下游 Benchmark 能力虽具有相关性，却并非可以直接互换的同一指标。

在实际训练中，算力和高质量语料均是有限资源。增大参数量和训练数据量会直接提高基础训练成本，数据清洗、筛选与重构会产生额外的质量提升成本，而长上下文训练还会引入随上下文窗口增长的注意力开销。领域配比虽然不必然增加 Token 总量，却决定有限数据预算在不同知识领域之间的分配。因此，模型性能提升实质上是参数规模、数据规模、数据质量、领域配比和上下文需求共同作用下的多资源配置问题。

基于此，本文围绕“数据属性刻画--广义损失建模--算力约束优化--能力前沿预测”展开研究：首先建立样本与领域质量评价体系并识别领域配比效应；其次将质量和配比因素引入经典标度律；随后在统一算力预算下求解参数、数据与质量投入的条件最优配置；最后连接预训练 Loss 与下游 Benchmark，分析历史能力提升来源并预测算力增长放缓情景下的未来能力前沿。

\subsection{具体问题}

\subsubsection{问题一重述}

利用附件 A 提供的多维质量信号，完成异构指标的统一量化，建立样本级和领域级训练数据质量评价体系，并说明不同质量维度发生冲突时的识别与综合方法。在此基础上，利用 $17$ 个训练领域的配比实验及其在 $13$ 个验证领域上的 Loss，刻画领域份额、领域组合与模型损失之间的关系，识别有利的调整方向和交互作用，并通过不同模型规模下的独立数据检验模型的稳定性与外推边界。

\subsubsection{问题二重述}

在经典参数量--数据量标度律的基础上，将问题一得到的数据质量与领域配比信息引入验证损失模型，建立关于参数量 $N$、训练 Token 数 $D$、质量水平 $Q$ 和领域配比 $\boldsymbol p$ 的广义标度律。利用附件 B 中不同来源和不同可信层级的数据完成模型选择、参数估计与外部检验，并据此比较四类因素对 Loss 的边际效用和弹性，讨论在给定性能目标下参数扩张、数据扩张与质量提升之间的有限替代关系，以及领域组合可能产生的互补效应。

\subsubsection{问题三重述}

在前两问模型的基础上，将基础训练、数据质量提升和长上下文注意力三类开销纳入统一算力预算，以最小化预测 Loss 为目标配置模型参数量、训练 Token 数和质量投入，并将最优 Token 总量分解到各训练领域。分别考察不同质量成本函数、上下文长度和算力预算对最优配置的影响，分析预算跨越不同数量级时各类成本份额、活跃约束与资源投入重点是否发生结构性转移，同时标记超出观测支持域的方案。

\subsubsection{问题四重述}

结合历史开源模型的训练规模、算力证据和多任务 Benchmark 记录，统一评测口径并构建综合能力指标；建立带有支持域约束的 Loss--Benchmark 映射，将前三问的资源配置结果转换到下游能力尺度。在此基础上，刻画开源模型能力前沿随训练算力和时间的演进，分解规模扩张与非规模技术因素对历史能力提升的相对贡献，并在算力增长延续、减速等情景下预测未来 $12$ 个月和 $24$ 个月的能力前沿及其不确定性区间。

 
